\newcommand{\pcell}[2]{\parbox[t]{#1}{\raggedright\strut#2\strut}}
\newcommand{\nsec}[1]{Ref.~\cite{AnalysisNote2026}, Sec.~\ref*{N-#1}}
\newcommand{\napp}[1]{Ref.~\cite{AnalysisNote2026}, App.~\ref*{N-#1}}

\section{Introduction}
\label{sec:intro}

Multidifferential neutrino cross sections test interaction models across muon and hadronic observables.
Their interpretation depends on correcting detector response and on assessing the bias and uncertainty of the
unfolding procedure.
Iterative Bayesian unfolding on a binned migration matrix~\cite{DAgostini:1994fjx} is widely used, but its
response tensor grows rapidly as observables are added.
OmniFold learns event weights in an unbinned feature space through classifier-based density-ratio
estimation~\cite{Andreassen:2019cjw,Andreassen:2021zzk}.
This makes simultaneous higher-dimensional unfolding possible without specifying a separate binned response tensor
for every projection.
Unbinned unfolding has been applied to collider data, for example to eight observables in deep-inelastic scattering
at H1~\cite{H1:2021wkz} and to twenty-four $Z$+jets observables at ATLAS~\cite{ATLAS:2024jets} (see
Ref.~\cite{Canelli:2025guide} for a survey).
In neutrino physics, it has been demonstrated on a public simulated T2K near-detector
sample~\cite{Huang:2025omnifold}.
Multidifferential neutrino measurements on data reach three
dimensions~\cite{MINERvA:2022qelike3d,MINERvA:2026qeleme,MicroBooNE:2023inc3d,NOvA:2026inc3d}.

We apply the method to MINERvA medium-energy $\nu_{\mu}$ data from the open-data release~\cite{MINERvA:OpenData}.
The result is a simultaneous charged-current inclusive cross section in five observables: the transverse and
longitudinal muon momenta $\pT$ and $\pPar$, the available hadronic energy $\Eavail$, the three-momentum transfer
$q_{3}$, and the hadronic invariant mass $W$.
It is reported in sixteen cells with its covariance (Sec.~\ref{sec:result16}) and compared with five generator
predictions (Sec.~\ref{sec:gen16}).
We first reproduce the published double-differential cross section
$d^{2}\sigma/(d\pT\,d\pPar)$~\cite{MINERvA:2021owq} as an end-to-end test of the unfolding, and we also give
three- and four-observable central values.
Reconstructed $W$ is fixed by the reconstructed $(\pT,\pPar,q_{3})$, so the five observables are simultaneous but
not independent.

The validation studies assess the result through the two-dimensional reproduction, background-subtraction tests,
recovery of generator-sized truth changes, a coverage test, and a fake-data study of estimator bias.
They examine whether the unfolding gives reliable multidimensional cross sections when the simulation mismodels the
truth by generator-sized amounts.
The coverage of unfolded intervals is the rate at which they contain the true cross section.

Interval coverage has not been validated.
The bias tests, the fake-data study, and the dependence of the generator $\chi^{2}$ on the flux correlations are
reported for the sixteen-cell result in Secs.~\ref{sec:fd16} and~\ref{sec:gen16}.

\emph{Scope.} The findings are those of one OmniFold implementation on one data set, for the observables,
configurations, and deformations tested, rather than of multidimensional OmniFold in general.
The limitations are listed in Sec.~\ref{sec:limitations}.

Section~\ref{sec:data} describes the data, the simulation and the observables, and Sec.~\ref{sec:method} the
unfolding and the reporting cells.
Section~\ref{sec:validation} gives the uncertainties and the validation studies, from the reproduction of the
published two-dimensional measurement to the fake-data tests of the sixteen-cell result, and
Sec.~\ref{sec:result16} gives the result and its comparison with generator predictions.
Section~\ref{sec:conclusions} discusses what limits the estimator, and Sec.~\ref{sec:limitations} summarizes the
principal limitations.
The companion analysis note~\cite{AnalysisNote2026} gives the full validation record, the construction of the
uncertainties, and the tests with their assumptions, and it is cited below by section.

\section{Data, simulation, and observables}
\label{sec:data}

The data come from the twelve MINERvA medium-energy, forward-horn-current run periods, corresponding to
$1.057\times10^{21}$ protons on target, with an average neutrino energy near \SI{6}{GeV}~\cite{MINERvA:2021owq}.
The input is an earlier production of the MINERvA open-data release (see Data availability), in the form of
pre-selected analysis ntuples (AnaTuples)~\cite{MINERvA:OpenData}.
The detector, including the downstream MINOS near detector that serves as a muon spectrometer, is described in
Ref.~\cite{MINERvA:2013zvz}.
The signal is a $\nu_{\mu}$ charged-current interaction, with any final state, whose true vertex lies in the
tracker fiducial volume ($5980<z<\SI{8422}{mm}$, inside a hexagon of apothem \SI{850}{mm}).
The signal muon has $\theta_{\mu}<20^{\circ}$, $1.5<\pPar<\SI{60}{GeV/c}$, and $\pT\le\SI{4.5}{GeV/c}$.
Selected simulated events outside this definition, including $\nu_{\mu}$ charged-current events whose muon lies
outside these limits, are background.
The selection requires a reconstructed vertex in the same volume and a reconstructed muon with
$\theta_{\mu}<20^{\circ}$ matched to a MINOS track of negative charge, together with a dead-time requirement, as in
Ref.~\cite{MINERvA:2021owq}.
The flux treatment follows Ref.~\cite{MINERvA:2016iqn}.
The analysis uses $4.12$ million selected data events.
The simulation contains $32.8$ million unweighted true signal events and $20.5$ million selected simulated events,
and it is scaled to the data exposure.
The simulated events are generated with GENIE~2.12.6~\cite{Andreopoulos:2009rq} and reweighted with the MINERvA
Tune~v1, which applies a non-resonant pion reduction, a random-phase-approximation correction, and an empirical
low-recoil two-particle--two-hole (2p2h) enhancement~\cite{MINERvA:2016lowrecoil}.
Event selection is held fixed as the classifier inputs and unfolded observables grow from $(\pT,\pPar)$ to
$(\pT,\pPar,\Eavail,q_3,W)$.
Adding features does not add selected events or relax the signal definition.

The unfolded observables are defined at truth level, with muon momentum components $\pT=p_{\mu}\sin\theta_{\mu}$
and $\pPar=p_{\mu}\cos\theta_{\mu}$.
The magnitude of the three-momentum transfer $q_{3}=|\vec{q}|$ comes from the generator's four-vectors, and the
hadronic invariant mass $W=\sqrt{M^{2}+2M(E_{\nu}-E_{\mu})-Q^{2}}$ uses the struck-nucleon mass $M$.
We define the available energy as
\begin{equation}
\Eavail=\sum_{p}(E_{p}-m_{p})+\sum_{\pi^{\pm}}(E_{\pi}-\SI{135}{MeV})
+\sum_{\pi^{0}}E_{\pi}+\sum_{\gamma}E_{\gamma},
\label{eq:eavail}
\end{equation}
which subtracts a defined constant of \SI{135}{MeV}, rather than the charged-pion mass, in the charged-pion term.
This constant is inherited from the analysis framework's definition.
It defines the reported observable and is used consistently by the analysis and every prediction.
The sums run over the generator's final-state particles after final-state interactions, where neutral pions appear
undecayed.
All other final-state particles, including neutrons, contribute nothing.
Reconstructed $\Eavail$ is $1.17$ times the calorimetric recoil energy in the tracker and electromagnetic
calorimeter, with muon deposits removed.
Reconstructed $q_{3}$ and $W$ instead combine the muon kinematics with a separate calorimetric estimate $q_{0}$ of
the energy transfer.

Reconstructed $W$ is an exact function of the reconstructed $(\pT,\pPar,q_{3})$, so the five reconstructed inputs
share information and do not provide five independent constraints.
The relation holds because $W_{\mathrm{reco}}^{2}=(\bar M+q_{0})^{2}-q_{3,\mathrm{reco}}^{2}$, with $\bar M$ an
average nucleon mass, and $q_{3,\mathrm{reco}}$ increases with $q_{0}$ at fixed muon kinematics.
This relation is exact apart from a clamp at $W^{2}=0$.
Only reconstructed $W$ is fixed by the other inputs.
Reconstructed $\Eavail$ comes from its own calorimetric estimator and is not a function of them.
At truth level, $\theta_{\mu}$, and with it $\pT$ and $\pPar$, is measured from a fixed beam axis rather than from
the true neutrino direction.
Truth $W$ is therefore fixed by $(\pT,\pPar,q_{3})$ only up to that event-by-event angle and the choice of $M$.
The definitions are detailed in Ref.~\cite{AnalysisNote2026}, Secs.~\ref*{N-sec:eavail-definition}
and~\ref*{N-sec:eavailw-definitions}.

The external predictions for CH are the untuned GENIE~2.12.10 without (``GENIE CV'') and with (``GENIE + MEC'') the
Valencia 2p2h component~\cite{Nieves:2011pp}, which models meson-exchange currents (MEC).
The predictions for carbon are NuWro~21.09~\cite{Golan:2012wx} and GiBUU~2019~\cite{Buss:2011mx}.
All predictions are per nucleon of their target.
The GENIE and NuWro predictions are reweighted exactly, event by event, to the analysis flux.
The GiBUU sample stops at $E_{\nu}\approx\SI{20}{GeV}$.
Unlike the analysis definitions, the predictions compute $W$ with the neutron mass for every event and measure
$\pT$ about their own neutrino axis.
Hydrogen has no quasielastic, 2p2h, or coherent $\nu_{\mu}$ charged-current channel, so normalizing per nucleon
does not remove the difference between the targets.
In GENIE's own CH sample, the carbon-only cross section per nucleon exceeds that of CH by \SI{8.3}{\percent} for
$W<\SI{1.1}{GeV}$.
The NuWro and GiBUU comparisons contain a target difference of this kind, which the sixteen-cell comparison of
Sec.~\ref{sec:gen16} does not correct.
The comparisons of the predictions with the analysis definitions are in Ref.~\cite{AnalysisNote2026},
Secs.~\ref*{N-sec:eavailw-axis} and~\ref*{N-sec:eavailw-generators}.

The fine grid is the analysis binning: the published $14\times16$ muon-momentum bins, crossed with 7 $\Eavail$,
7 $q_{3}$, and 6 $W$ bins.
A coarse $3^{5}$ grid in $(\pT,\pPar,\Eavail,q_{3},W)$ has its edges at the fine edges nearest the simulated
tertiles.
A coarse cell is kept if its expected reconstructed signal and its truth weight at the data exposure are both at
least 1000 events.
We call this requirement the support rule.
Of the coarse cells, \jtNcells{} pass; these are the joint cells.
The others, which hold \SI{0.47}{\percent} of the expected signal, stay in the total only.
The construction of the coarse grid is given in \nsec{sec:joint5d-design}.

\section{Cross-section extraction}
\label{sec:method}

Each OmniFold iteration has two classification steps.
A reconstruction-level classifier distinguishes data from simulated events and defines a likelihood-ratio weight.
The event pairing transfers that weight to truth, including unmatched events, and a truth-level classifier
propagates a smooth reweighting to the next iteration.
Cross sections are formed by histogramming the reweighted truth events in the reporting binning and dividing by
integrated flux, target count, and bin width.
Events that fail reconstruction enter as unmatched truth events, so no separate efficiency correction is applied.

Backgrounds enter the reconstruction-level step through the density the classifier learns to match.
In the purity-based subtraction, every data event in reconstructed bin $b$ carries the weight
$w_{\mathrm{pur}}(b)=\max(0,\,1-B_{b}/D_{b})$, where $D_{b}$ and $B_{b}$ are the data and predicted background
counts in that bin, so the classifier's target density is $D(x)\,w_{\mathrm{pur}}(b)$.
The unbinned negative-weight treatment instead adds simulated background events with negative weight to the data,
so that the target is $D(x)-B(x)$ pointwise.
A classifier refinement~\cite{StayPositive2025}, clipped at zero, converts the signed sample to non-negative weights
before the first iteration.
The two targets share bin integrals wherever $B_{b}<D_{b}$ but differ inside a bin wherever the background fraction
varies within it (\napp{app:negweight}).

Table~\ref{tab:estimators} lists the two estimators and the results each is used for.
The two-observable central value uses LightGBM at seed~42~\cite{Ke:2017lightgbm} with negative-weight background subtraction and a
refinement of 400 trees and 31 leaves.
This reproducible central value is that of its own systematic-universe ensemble, so its covariance is built with
the same estimator.
The three- and four-observable central values and the sixteen-cell five-observable result (Sec.~\ref{sec:result16})
use the joint-test estimator, a deterministic LightGBM with the same background treatment and refinement capacity.
This estimator removes most of the bias of purity-based subtraction at nominal truth in the highest-$W$ cells
(Sec.~\ref{sec:validation}), and the note's simulation-calibrated tests use the same estimator.
The variants used in development, and the results that rest on them, are described in \nsec{sec:method-config}.

Both estimators use five OmniFold iterations with gradient-boosted decision-tree
classifiers~\cite{Pedregosa:2011sklearn,Ke:2017lightgbm}.
The iteration count and classifier capacity regularize the unfolding.
The nominal settings are the scikit-learn defaults, with LightGBM matched to their capacity.
The refinement capacity of the joint-test estimator was chosen in simulation closure studies.
No hyperparameter search was performed on data.
Each classifier is trained on the full sample whose weights it then sets, without a held-out or cross-fitted split,
early stopping, or class balancing (\nsec{sec:method-algo}).

The joint-test estimator keeps five iterations because 10, 15, or 20 iterations did not reduce its model
dependence, measured as the largest residual proxy over the truth changes used in its development.
The comparison and its values are given in \nsec{sec:gbdt-model-dependence}.

\begin{table}[tb]
\caption{Estimators. Both use five OmniFold iterations, gradient-boosted trees of nominal capacity (100 trees,
8 leaves) in both OmniFold steps, and refined negative-weight background subtraction. The joint-test estimator is
a deterministic LightGBM (fixed seeds, reproducible training), and the refinement classifier of the two-observable
estimator is seeded at the estimator seed plus three. Results from the two estimators are not combined.
\nsec{sec:method-config} gives the full matrix, including the variants used in development.}
\label{tab:estimators}
\begin{tabular}{@{}lll@{}} \toprule
\pcell{0.24\columnwidth}{Estimator} & \pcell{0.30\columnwidth}{Background treatment} & \pcell{0.38\columnwidth}{Used for} \\ \midrule
\pcell{0.24\columnwidth}{joint-test estimator} & \pcell{0.30\columnwidth}{negative weights; refinement 400 trees, 31 leaves} & \pcell{0.38\columnwidth}{3D and 4D central values; sixteen-cell 5D result (Sec.~\ref{sec:result16}); held-out shapes (Fig.~\ref{fig:deformation}b); precision study} \\
\pcell{0.24\columnwidth}{two-observable estimator: LightGBM, seed 42} & \pcell{0.30\columnwidth}{negative weights; refinement 400 trees, 31 leaves} & \pcell{0.38\columnwidth}{2D central value and covariance (Sec.~\ref{sec:repro2d})} \\ \bottomrule
\end{tabular}
\end{table}

\subsection{Reporting definitions and the precision target}
\label{sec:uncertainty-status}

To include model dependence, a dedicated study defines the envelope of a cell as the largest shift of its data
result when the simulation's prior is reweighted to each of five generator-sized changes of the truth, and adds it
linearly to the probabilistic half-width.
The envelope is a sensitivity diagnostic that tracks the regularization bias, not an interval.
Our precision target is a \SI{10}{\percent} median half-width for the total 68\% interval.
With the envelope added, none of the reporting definitions of a five-observable measurement tested in the study
meets this target.
With the joint-test estimator, the envelope alone exceeds it for the \jtNcells{} coarse joint cells and for the
fallback definition, the 39 supported $(\Eavail,W)$ cells.
For a coarser, non-nested partition of all five observables into 27 cells, the same estimator's envelope alone does not exceed it,
but the \SI{1.4}{\percent} normalization contained in every probabilistic half-width lifts the median above it.
Over these definitions, the envelope has a median of about 10--13\% per cell, comparable to the generator
differences that such a measurement would need to resolve.
The study is described in \nsec{sec:gbdt-model-dependence}.

\subsection{Reporting cells}

A five-observable cross section with total uncertainties is reported separately from those reporting definitions:
the joint-test estimator's data result on 16 cells, R00--R15.
They come from the \fvval{h2-nboxes} boxes formed by one interior edge per axis, at $\pT$ \fvval{h2-mid-pt},
$\pPar$ \fvval{h2-mid-pz} and $q_3$ \fvval{h2-mid-q3}~GeV/$c$, and at $\Eavail$ \fvval{h2-mid-eavail} and
$W$ \fvval{h2-mid-W}~GeV.
In forming the cells, each supported box with purity below \fvval{h2-tau} is merged into a neighbour,
\fvval{h2-nunreported} boxes remain unreported, and every reported cell passes the support rule.
The result's covariance is given in Sec.~\ref{sec:cov16}.
